% ABSTRACT -- fill every \TODO after running fit.py. Keep to ~150-200 words.
Hoffmann et al.\ (2022) fixed the field's default recipe for allocating a training compute budget between model size and training tokens (``Chinchilla'', roughly 20 tokens per parameter). Schaeffer et al.\ (2025) recently stress-tested this recipe: they found that the conclusion is insensitive to how Chinchilla's model sizes are computed, but can be affected by additive or size-dependent errors in the parameter counts. We carry out a small-scale analogue of this analysis on a single consumer GPU. We train \TODO{15} GPT-style language models (0.39M--59M non-embedding parameters) on C4 at \TODO{three} FLOP budgets ($3\times10^{15}$--$3\times10^{16}$) and fit all three of Hoffmann et al.'s estimation approaches (training-curve envelope, IsoFLOP parabolas, parametric loss surface) under three parameter-count conventions (total, non-embedding, non-embedding plus output head), and apply four structured perturbations to the parameter counts following Schaeffer et al. \TODO{One sentence: the main finding about how the fitted optimal token count $D^*$ and exponent $a$ change across conventions.} \TODO{One sentence: how the three approaches compare.} A grid-shrinking sweep shows that \TODO{one sentence: when the conclusions become unstable as the fitting grid shrinks.} Code is released to support replication.
